\documentclass[11pt]{article}
\usepackage{mathunal-base}
\mathunalfuente{serif}

\renewcommand{\examtitulo}{Primer Examen Parcial de Álgebra Lineal}
\renewcommand{\examfecha}{15 de Septimbre de 2014}

\newcommand{\vf}{\fbox{V}\ \fbox{F}}

\begin{document}

\examheader

\vspace{4pt}
\noindent
\begin{minipage}[t]{0.62\linewidth}
\textbf{Puntaje. Sólo para uso Oficial}\\[4pt]
\begin{tabular}{|c|c|c|c|c|c|}
\hline
1--5 & 6 & 7 & 8 & \textbf{TOTAL} & \textbf{NOTA} \\ \hline
 & & & & & \\ \hline
\end{tabular}
\end{minipage}

\vspace{6pt}
\noindent\textbf{Instrucciones:} La duración del examen es de 1 hora y 40 minutos. El examen consta de ocho preguntas en tres hojas impresas por ambos lados, verifique que su examen esté completo. En las preguntas con procedimiento justifique sus respuestas en los espacios asignados. No está permitido el uso de calculadoras, sacar hojas en blanco ni ningún tipo de apuntes durante el examen. Tampoco se permiten preguntas. Por favor verifique que su celular esté apagado.

\vspace{6pt}
\noindent\textbf{IDENTIFICACIÓN}

\vspace{8pt}
\noindent Nombre: \dotfill\ Cédula \dotfill

\vspace{4pt}
\noindent Profesor: \dotfill\ Grupo \dotfill

\vspace{10pt}
\begin{center}\textbf{I. Opción Múltiple y Completación}\end{center}
\noindent En las preguntas 1 a 8 complete los espacios en blanco o elija la opción correcta según el caso. Marque con una X la letra de su elección. \textbf{NOTA}: En esta sección se califica sólo la respuesta y no se tiene en cuenta el procedimiento.

\begin{enumerate}
\item{}[7pt] Considere los vectores $\mathbf{u} = \begin{bmatrix}0\\1\\1\end{bmatrix}$ y $\mathbf{v} = \begin{bmatrix}1\\0\\1\end{bmatrix}$. El ángulo entre los vectores $\mathbf{u}$ y $\mathbf{v}$ es

\vspace{4pt}
\fbox{\begin{minipage}[t][2.3cm]{\dimexpr\linewidth-2\fboxsep-2\fboxrule}
$\theta = $
\end{minipage}}

\item{}[12pt] Indique con una X si los siguientes enunciados son verdaderos o falsos. Cada numeral vale 3 puntos.

\vspace{6pt}
\noindent (a) Todo sistema de ecuaciones lineales con $n$ variables y $n$ ecuaciones tiene una única solución.

\noindent (V)\ \dotfill\ (F).\ \dotfill

\vspace{6pt}
\noindent (b) Si $\mathbf{v_1}, \mathbf{v_2}, \mathbf{v_3}$ y $\mathbf{v_4}$ son vectores en $\mathbb{R}^3$ entonces $\{\mathbf{v_1},\mathbf{v_2},\mathbf{v_3},\mathbf{v_4}\}$ es un conjunto linealmente dependiente.

\noindent (V)\ \dotfill\ (F).\ \dotfill

\vspace{6pt}
\noindent (c) Todo sistema de ecuaciones lineales con más variables que ecuaciones tiene al menos una solución.

\noindent (V)\ \dotfill\ (F).\ \dotfill

\vspace{6pt}
\noindent (d) Si $\mathbf{v_1}, \mathbf{v_2}$ y $\mathbf{v_3}$ son vectores linealmente dependientes entonces $\mathbf{v_1}$ es un múltiplo de $\mathbf{v_2}$ o de $\mathbf{v_3}$.

\noindent (V)\ \dotfill\ (F).\ \dotfill

\item{}[7pt] Sea $A$ la matriz $A = \begin{bmatrix}1&-2&-2\\1&1&1\\2&-1&-1\end{bmatrix}$. Entonces

\vspace{4pt}
\fbox{\begin{minipage}[t][2.3cm]{\dimexpr\linewidth-2\fboxsep-2\fboxrule}
$\mathrm{Rango}(A) = $
\end{minipage}}

\item{}[7pt] Considere los vectores $\mathbf{i} = \begin{bmatrix}1\\0\\0\end{bmatrix}$, $\mathbf{j} = \begin{bmatrix}0\\1\\0\end{bmatrix}$ y $\mathbf{x} = \begin{bmatrix}1\\1\\2\end{bmatrix}$. ¿Cuáles de las siguientes afirmaciones son verdaderas?

\vspace{4pt}
\noindent $I.\ \mathbf{i}-\mathbf{j}$ es ortogonal a $\mathbf{x}$ \hspace{1cm} $II.\ \mathbf{x}-\mathbf{i}$ es paralelo a $\mathbf{j}$ \hspace{1cm} $III.\ \mathbf{i}+\mathbf{j}$ no es ortogonal a $\mathbf{x}$

\vspace{6pt}
\noindent
\begin{minipage}[t]{0.33\linewidth}
$(a)$ $I$ y $II$ únicamente
\end{minipage}%
\begin{minipage}[t]{0.33\linewidth}
$(b)$ $II$ únicamente
\end{minipage}%
\begin{minipage}[t]{0.33\linewidth}
$(c)$ $III$ únicamente
\end{minipage}

\vspace{4pt}
\noindent
\begin{minipage}[t]{0.5\linewidth}
$(d)$ $I$ y $III$ únicamente
\end{minipage}%
\begin{minipage}[t]{0.5\linewidth}
$(e)$ Ninguna de las anteriores
\end{minipage}

\item{}[7pt] Considere el siguiente sistema de ecuaciones lineales
\[
\left\{
\begin{aligned}
x+y+z &=1\\
x-z &=2
\end{aligned}
\right.
\]
¿Cuál de las siguientes afirmaciones es verdadera?

\vspace{6pt}
\noindent
\begin{minipage}[t]{0.5\linewidth}
$(a)$ El sistema no tiene solución,
\end{minipage}%
\begin{minipage}[t]{0.46\linewidth}
$(b)$ El sistema tiene infinitas soluciones,
\end{minipage}

\vspace{4pt}
\noindent
\begin{minipage}[t]{0.5\linewidth}
$(c)$ El sistema tiene solución única,
\end{minipage}%
\begin{minipage}[t]{0.46\linewidth}
$(d)$ El sistema tiene exactamente dos soluciones.
\end{minipage}
\end{enumerate}

\vspace{6pt}
\begin{center}\textbf{II. Solución con Procedimiento}\end{center}

\begin{enumerate}
\setcounter{enumi}{5}
\item{}[13pt] Sean $\mathbf{u}, \mathbf{v}$ y $\mathbf{w}$ vectores en $\mathbb{R}^n$. Demostrar que $\mathbf{u}, \mathbf{v}$ y $\mathbf{w}$ pertenecen a $gen\{\mathbf{u}, \mathbf{u}-\mathbf{v}, \mathbf{u}+\mathbf{v}-\mathbf{w}\}$.

\vspace{6cm}

\item{}[25pt] Consideremos la siguiente red de transporte

\vspace{6pt}
\begin{center}
\begin{tikzpicture}[>=Stealth, scale=1]
  \coordinate (A) at (0,2);
  \coordinate (B) at (0,-2);
  \coordinate (C) at (3.6,0);
  \draw[thick] (0,3.2) -- (0,-3.2);
  \draw[thick,->] (A) -- node[above,sloped]{$x_2$} (C);
  \draw[thick,->] (A) -- node[left]{$x_1$} (B);
  \draw[thick,->] (B) -- node[below,sloped]{$x_3$} (C);
  \fill (A) circle (2pt) node[above right]{A};
  \fill (B) circle (2pt) node[below right]{B};
  \fill (C) circle (2pt) node[right]{C};
  \draw[->] (0,3.2) -- (0,2.35) node[midway,right]{15};
  \draw[->] (-1.3,3.1) -- (A) node[midway,above left]{10};
  \draw[->] (-1.3,-1.3) -- (B) node[midway,above left]{5};
  \draw[->] (B) -- (0,-3.2) node[midway,right]{20};
  \draw[->] (C) -- (4.9,1.1) node[midway,above]{5};
  \draw[->] (C) -- (4.9,-1.1) node[midway,below]{5};
\end{tikzpicture}
\end{center}

\noindent Responda las preguntas planteadas en los numerales (a), (b) y (c).

\begin{enumerate}
\item{}[12pt] Establezca y resuelva un sistema de ecuaciones lineales para encontrar los flujos posibles en la red de la figura anterior.

\vspace{5cm}

\item{}[7pt] Si las direcciones de los flujos se respetan, ¿cuáles son los posibles flujos mínimos y máximos a través de cada rama?

\vspace{4cm}

\item{}[6pt] Suponiendo que las direcciones de los flujos se respetan, halle los valores de $x_1, x_2$ y $x_3$ de tal forma que se obtenga el mínimo flujo vehicular en la rama $\overline{AB}$.
\end{enumerate}

\item{}[22pt]
\begin{enumerate}
\item{}[10pt] Sean $\mathbf{v_1} = \begin{bmatrix}1\\1\\0\end{bmatrix}$, $\mathbf{v_2} = \begin{bmatrix}-1\\0\\1\end{bmatrix}$ y $\mathbf{b} = \begin{bmatrix}1\\2\\3\end{bmatrix}$. Determine si $\mathbf{b}$ es combinación lineal de los vectores $\mathbf{v_1}$ y $\mathbf{v_2}$.

\vspace{6cm}

\item{}[12pt] Determine si los vectores $\mathbf{w_1} = \begin{bmatrix}1\\0\\1\\0\end{bmatrix}$, $\mathbf{w_2} = \begin{bmatrix}-1\\0\\0\\1\end{bmatrix}$ y $\mathbf{w_3} = \begin{bmatrix}1\\1\\1\\-1\end{bmatrix}$ son linealmente independientes o linealmente dependientes.

\vspace{6cm}
\end{enumerate}
\end{enumerate}

\examfooter
\end{document}
